\chapter{Transformations and Iteration Axes}\label{ch:transformations}

After stage~1 the kernel body is a typed block of assignments, without any loops. This
chapter explains how the iteration strategy is attached (stage~2), which optimizations
follow (stage~3), and how the IR is lowered to plain C plus intrinsics (stage~4). It
ends by running the whole pipeline by hand. All passes live in
\file{backend/transformations/} and are called as \code{Pass(ctx)(ast)}.

\section{Iteration axes}\label{sec:axes}

\ps{} does not generate loops directly. It first represents the iteration space as an
\term{axes cube}: a \cls{PsAxesCube} node with one \cls{PsAxisRange}
\code{(counter, start, stop, step)} per dimension, in loop order, around the kernel body.
An \term{expansion strategy} then takes the dimensions off the cube, one at a time from
the outermost, and wraps the body in \term{iteration axes}:
\begin{center}\small
\begin{tabularx}{\linewidth}{@{}l l >{\raggedright\arraybackslash}X@{}}
  \toprule
  \textbf{Expansion} (\cls{AxisExpansion}) & \textbf{Axis node} & \textbf{Meaning}\\
  \midrule
  \code{loop()} & \cls{PsLoopAxis} & sequential loop\\
  \code{parallel_loop(num_threads, schedule, collapse)} & \cls{PsParallelLoopAxis} & OpenMP work-sharing loop\\
  \code{block_loop(k)}, \code{parallel_block_loop(k)} & loop axes & loop over blocks of $k$ iterations (loop blocking, SIMD blocks)\\
  \code{peel_for_divisibility(k)} & -- & split the dimension into a part divisible by $k$ and a remainder; the next strategy entries are two lists, one per part\\
  \code{simd(lanes)} & \cls{PsSimdAxis} & execute \code{lanes} iterations in one vector operation\\
  \code{gpu_block}, \code{gpu_thread}, \code{gpu_block_x_thread} & GPU axes & map the dimension to \code{blockIdx}, \code{threadIdx} or both\\
  \code{gridstrided_loop()} & loop axis & loop with the GPU grid size as stride\\
  \bottomrule
\end{tabularx}
\end{center}
A strategy is written as a list, from the outermost to the innermost dimension, and
compiled with \code{ae.create_strategy([...])}. The CPU strategy of the driver
(\cls{DefaultCpuLoopStrategies}, \file{codegen/cpu_loop_strategies.py}) builds the
following shapes, with \code{L} the number of vector lanes:
\begin{plaincode}
not vectorized:  [parallel_loop | loop, loop, ...]
vectorized:      [parallel_loop | loop, loop, ...,
                  peel_for_divisibility(L),
                  [block_loop(L, assume_divisible=True), simd(L)],   # vector part
                  [loop()]]                                          # remainder
loop blocking:   [block loops for the blocked dimensions] + one of the above
\end{plaincode}
For a 1D kernel the block loop of the vector part itself becomes the parallel loop.
Before the expansion, the CPU strategy applies \code{assume_inner_stride_one} by
replacing the innermost buffer strides with the constant 1. Afterwards it runs
\cls{HoistIterationInvariantDeclarations}. On the axes tree this moves declarations out
of every axis whose \code{can_hoist} flag is set.

The GPU strategy (\cls{GpuIndexMappingStrategy}, \file{codegen/gpu_indexing.py}) maps
the indexing schemes of \cref{sec:gpu-options} onto these expansions:
\code{Linear3D} uses \code{gpu_block_x_thread} for every dimension,
\code{GridstridedLinear3D} wraps grid-strided loops around them, and
\code{Blockwise4D} uses \code{gpu_block} for the slow dimensions and
\code{gpu_thread("x")} for the fastest one.

\section{Materializing the axes}

\cls{MaterializeAxes} (\file{materialize_axes.py}) replaces every axis by real code:
\begin{description}
  \item[\cls{PsLoopAxis}] becomes a \cls{PsLoop}.
  \item[\cls{PsParallelLoopAxis}] becomes
    \code{#pragma omp parallel} (with \code{num_threads} and \code{reduction} clauses)
    around a block that contains \code{#pragma omp for} (with \code{schedule} and
    \code{collapse}) and the loop.
  \item[\cls{PsSimdAxis}] is handed to the \cls{AstVectorizer}. Reduction accumulators
    get a vector copy that is initialized by broadcasting the neutral element and is
    combined into the scalar at the end with \cls{PsVecHorizontal}.
  \item[GPU axes] become declarations of the counters from \code{blockIdx},
    \code{blockDim} and \code{threadIdx}. All bounds checks are merged into one guard
    \code{if}.
\end{description}

\paragraph{The vectorizer.}
\cls{AstVectorizer} (\file{ast_vectorizer.py}) builds a SIMD copy of a scalar subtree.
Memory accesses whose indices do not depend on the vectorized counter are
\emph{lane-invariant} and are broadcast (\cls{PsVecBroadcast}). Accesses whose index is
an affine function of the counter become \cls{PsVecMemAcc}, either contiguous or with a
stride. Every symbol declared inside the subtree gets a vectorized copy. Anything else
raises a \code{VectorizationError}. The vectorizer performs \emph{no legality checks}:
it relies on the kernel analysis having ruled out loop-carried dependencies.

\paragraph{Reductions.}
\cls{ReductionsToMemory} completes the reduction pattern. It prepends
\code{<sym>_local = <neutral element>} to the kernel and appends a write-back
\code{<sym>[0] = op(<sym>[0], <sym>_local)}, which the platform lowers later
(\code{resolve_reduction}: a plain update on CPUs, atomics on CUDA). Compare
\cref{lst:reduction}.

\section{Optimizations}

\begin{description}
  \item[\cls{CanonicalizeSymbols}] establishes the canonical form (\cref{sec:ast}). It
    renames symbols that are declared more than once and makes every symbol that is never
    updated \code{const}. It runs on the axes cube, after the iteration space is
    materialized, and at the very end.
  \item[\cls{EliminateConstants}] folds constant integer (and relational) expressions,
    removes identities ($x\cdot 1$, $x+0$, $x\cdot 0$), and with
    \code{extract_constant_exprs=True} moves constant subexpressions such as
    $1.0/4.0$ into declarations \code{__c_...} at the top of the kernel. Floating-point
    folding is off by default (\code{fold_floats=False}), so $1.0/4.0$ is still
    computed, once, by the C code.
  \item[\cls{HoistIterationInvariantDeclarations}] moves declarations whose value does
    not depend on a loop counter (or on anything written in the loop) out of loops and
    hoistable axes. This is how \code{ctr_1__rem_start} and the broadcast constants in
    \cref{lst:avx} end up outside the loop nest.
\end{description}

\section{Lowering}

\begin{description}
  \item[\cls{SelectIntrinsics}] (vector CPU platforms only). Each vector platform
    provides a subclass via \code{platform.get_intrinsic_selector()}. It maps vector types
    to intrinsic types and vector operations, broadcasts, loads, stores, casts,
    horizontal reductions and math functions to \code{CFunction} calls.
    \Cref{sec:vector-platforms} describes the interface.
  \item[\cls{LowerToC}] linearizes every \cls{PsBufferAcc} into a \cls{PsMemAcc} on the
    base pointer, with index $\sum_i \text{idx}_i\cdot\text{stride}_i$. Anonymous struct
    types are replaced by byte offsets on a \code{uint8_t} pointer (see the sparse kernel
    in \cref{sec:sparse}).
  \item[\cls{SelectFunctions}] replaces every remaining IR function call
    (\code{exp}, \code{min}, $\pi$, GPU indices, RNG, reduction write-back) with the
    platform's implementation (\code{platform.select_function(call)} and
    \code{platform.resolve_reduction(...)}).
\end{description}

\section{Overview of all passes}

\begin{center}\small
\begin{tabularx}{\linewidth}{@{}l l >{\raggedright\arraybackslash}X@{}}
  \toprule
  \textbf{Pass} & \textbf{Module} & \textbf{Used by}\\
  \midrule
  \cls{CanonicalizeSymbols} & \file{canonicalize_symbols.py} & driver (stages 3 and 5), axis strategies\\
  \cls{AxisExpansion} & \file{axis_expansion.py} & CPU and GPU axis strategies\\
  \cls{HoistIterationInvariantDeclarations} & \file{hoist_iteration_invariant_decls.py} & axis strategies, stage 3\\
  \cls{MaterializeAxes} & \file{materialize_axes.py} & stage 2\\
  \cls{AstVectorizer} & \file{ast_vectorizer.py} & \cls{MaterializeAxes}\\
  \cls{ReductionsToMemory} & \file{reductions_to_memory.py} & stage 2\\
  \cls{EliminateConstants} & \file{eliminate_constants.py} & stage 3, many helpers\\
  \cls{SelectIntrinsics} & \file{select_intrinsics.py} & stage 4 (vector CPUs)\\
  \cls{LowerToC} & \file{lower_to_c.py} & stage 4, vectorizer\\
  \cls{SelectFunctions} & \file{select_functions.py} & stage 4\\
  \cls{CanonicalClone} & \file{canonical_clone.py} & peeling, loop reshaping\\
  \midrule
  \cls{EliminateBranches} & \file{eliminate_branches.py} & \emph{not used by the driver} (optional islpy analysis)\\
  \cls{ReshapeLoops} & \file{reshape_loops.py} & \emph{not used by the driver} (peel, cut loops)\\
  \cls{AddOpenMP}, \cls{InsertPragmasAtLoops} & \file{add_pragmas.py} & \emph{not used by the driver}; superseded by parallel loop axes\\
  \bottomrule
\end{tabularx}
\end{center}

\section{Running the pipeline by hand}\label{sec:manual-pipeline}

Because every stage is a callable object, the pipeline can be executed step by step.
This is the best way to understand it, and the starting point for experiments with new
passes. The following example is modelled on the upstream walkthrough
(\file{docs/source/backend/backend_walkthrough.md}). It generates an OpenMP-parallel AVX
kernel for $g = c\,f + 1$ on \code{fzyx} fields.

\paragraph{1. Context, iteration space and parsing.}
\pysnippet{ex11_pipeline.kernel}
\pysnippet{ex11_pipeline.context}
\irsnippet[nolol=false,caption={After parsing: a typed block without loops.}]{ex11_pipeline.s1-parsed}
The driver's CPU strategy applies the unit-stride assumption at this point. Here it is
done by hand:
\pysnippet{ex11_pipeline.unit-stride}

\paragraph{2. The axes cube.}
\pysnippet{ex11_pipeline.axes}
\irsnippet[nolol=false,caption={The iteration space as an axes cube (slowest dimension first).}]{ex11_pipeline.s2-cube}

\paragraph{3. Axis expansion.}
\pysnippet{ex11_pipeline.expand}
\irsnippet[nolol=false,caption={The axes tree: a parallel loop over \texttt{ctr\_1}; \texttt{ctr\_0} split into SIMD blocks and a remainder.}]{ex11_pipeline.s3-axes}

\paragraph{4. Materialization and optimization.}
\pysnippet{ex11_pipeline.materialize}
\irsnippet[nolol=false,caption={Loops, OpenMP pragmas and vector nodes (\texttt{vec\_memacc}, \texttt{vec\_broadcast}).}]{ex11_pipeline.s4-materialized}

\paragraph{5. Lowering for x86 AVX.}
\pysnippet{ex11_pipeline.lowering}
\irsnippet[nolol=false,caption={After intrinsic selection, \texttt{LowerToC} and \texttt{SelectFunctions}.}]{ex11_pipeline.s5-lowered}

\paragraph{6. Wrapping up.}
\pysnippet{ex11_pipeline.finalize}
\csnippet[nolol=false,caption={The hand-made kernel.}]{ex11_pipeline.final}
The result is almost identical to what \code{create_kernel} produces for the equivalent
configuration (\file{generated/ex11_pipeline.driver.c}). There are two differences:
\code{_data_f} is not \code{const}, because \cls{KernelAnalysis} did not run and so did
not mark the read-only field, and the \code{omp for} pragma lacks \code{schedule(static)},
which the driver takes from \code{cpu.openmp.schedule}.
